% Options for packages loaded elsewhere
\PassOptionsToPackage{unicode}{hyperref}
\PassOptionsToPackage{hyphens}{url}
\documentclass[
]{article}
\usepackage{xcolor}
\usepackage{amsmath,amssymb}
\setcounter{secnumdepth}{-\maxdimen} % remove section numbering
\usepackage{iftex}
\ifPDFTeX
  \usepackage[T1]{fontenc}
  \usepackage[utf8]{inputenc}
  \usepackage{textcomp} % provide euro and other symbols
\else % if luatex or xetex
  \usepackage{unicode-math} % this also loads fontspec
  \defaultfontfeatures{Scale=MatchLowercase}
  \defaultfontfeatures[\rmfamily]{Ligatures=TeX,Scale=1}
\fi
\usepackage{lmodern}
\ifPDFTeX\else
  % xetex/luatex font selection
\fi
% Use upquote if available, for straight quotes in verbatim environments
\IfFileExists{upquote.sty}{\usepackage{upquote}}{}
\IfFileExists{microtype.sty}{% use microtype if available
  \usepackage[]{microtype}
  \UseMicrotypeSet[protrusion]{basicmath} % disable protrusion for tt fonts
}{}
\makeatletter
\@ifundefined{KOMAClassName}{% if non-KOMA class
  \IfFileExists{parskip.sty}{%
    \usepackage{parskip}
  }{% else
    \setlength{\parindent}{0pt}
    \setlength{\parskip}{6pt plus 2pt minus 1pt}}
}{% if KOMA class
  \KOMAoptions{parskip=half}}
\makeatother
\usepackage{longtable,booktabs,array}
\usepackage{calc} % for calculating minipage widths
% Correct order of tables after \paragraph or \subparagraph
\usepackage{etoolbox}
\makeatletter
\patchcmd\longtable{\par}{\if@noskipsec\mbox{}\fi\par}{}{}
\makeatother
% Allow footnotes in longtable head/foot
\IfFileExists{footnotehyper.sty}{\usepackage{footnotehyper}}{\usepackage{footnote}}
\makesavenoteenv{longtable}
\setlength{\emergencystretch}{3em} % prevent overfull lines
\providecommand{\tightlist}{%
  \setlength{\itemsep}{0pt}\setlength{\parskip}{0pt}}
\usepackage{bookmark}
\IfFileExists{xurl.sty}{\usepackage{xurl}}{} % add URL line breaks if available
\urlstyle{same}
\hypersetup{
  pdftitle={Uncertainty-aware Multi-view Expert Fusion for Interpretable Candlestick Image Trend Classification},
  hidelinks,
  pdfcreator={LaTeX via pandoc}}

\title{Uncertainty-aware Multi-view Expert Fusion for Interpretable
Candlestick Image Trend Classification}
\author{}
\date{}

\begin{document}
\maketitle

\section{Uncertainty-aware Multi-view Expert Fusion for Interpretable
Candlestick Image Trend
Classification}\label{uncertainty-aware-multi-view-expert-fusion-for-interpretable-candlestick-image-trend-classification}

\subsection{Abstract}\label{abstract}

Candlestick-chart images provide a visually intuitive representation of
financial time series and have been increasingly used for stock trend
classification. However, high accuracy in chart-image classification is
often difficult to interpret: a model may recognize the trend already
visible inside the input window rather than predict the subsequent price
movement, and standard feature-fusion models provide limited insight
into which technical views support a decision. This paper studies
candlestick image trend classification from the perspective of
interpretable multi-expert prediction. We propose an
\textbf{Uncertainty-aware Multi-view Expert Fusion} framework, termed
\textbf{U-MoE}, that treats different technical chart renderings as
view-specific image experts. Each stock window is rendered into four
pure image views: raw candlestick, moving-average candlestick, RSI, and
MACD. A shared ResNet18 encoder extracts view-level visual
representations, each expert produces a class distribution, and a gating
module dynamically fuses expert predictions. Crucially, the gate is
regularized by expert predictive entropy, allowing uncertain views to
receive lower decision weight.

We evaluate U-MoE under strict chronological splits on 5,283 samples
constructed from ten U.S. stocks. Our experiments show a clear
distinction between visible trend recognition and future direction
prediction. A reproduced paper-style window-trend classifier reaches
98.37\% test accuracy, whereas strict next-day prediction remains near
random. On the window-trend task, black-box concat fusion achieves the
highest raw accuracy of 98.12\%, while U-MoE achieves 96.74\% accuracy
with explicit view-level expert weights and uncertainty estimates. On
next-day prediction, U-MoE does not improve accuracy, but its expert
entropies approach the maximum binary entropy, revealing that the model
cannot identify reliable visual evidence for the future label. These
results support a cautious but useful conclusion: candlestick images are
effective for recognizing visible trend structures, while U-MoE provides
an interpretable mechanism for diagnosing which chart views contribute
to a prediction and when the model should be considered uncertain.

\textbf{Keywords:} candlestick chart, financial image analysis, mixture
of experts, uncertainty estimation, interpretable machine learning,
multi-view learning, Grad-CAM

\subsection{1. Introduction}\label{introduction}

Financial time-series prediction is a challenging problem because asset
prices are noisy, non-stationary, and shaped by latent market mechanisms
that are only partially observable. In recent years, an alternative to
numerical feature engineering has gained attention: converting
historical price windows into candlestick-chart images and applying
computer vision models to classify future or contemporaneous market
trends. This line of work is appealing because candlestick charts encode
open, high, low, and close prices in a compact visual format familiar to
market participants, and convolutional neural networks can learn local
visual patterns from such images.

Despite promising reported results, image-based stock prediction raises
two important questions. First, what exactly is being predicted? In many
chart-image settings, the label describes whether the displayed window
itself moves upward or downward. A model can then obtain high accuracy
by recognizing a trend already visible inside the input image. This is a
legitimate visual trend recognition task, but it differs from predicting
the next trading day's price direction. Second, when a model uses
multiple technical chart views, such as candlesticks, moving averages,
RSI, and MACD, how can we know which view contributes to the final
decision? Standard feature concatenation can be accurate, but it
provides little view-level interpretability.

This paper focuses on the second issue while explicitly controlling for
the first. We aim to build an interpretable pure-image model for
candlestick trend classification, not to overstate the predictability of
short-term returns. Our central hypothesis is that different technical
image views behave like distinct experts: a candlestick view may capture
price body and wick patterns, a moving-average view may capture trend
direction, RSI may capture momentum saturation, and MACD may capture
oscillator changes. A mixture-of-experts architecture is therefore a
natural fit for multi-view financial chart images, provided that the
model can also quantify when an expert is uncertain.

We propose \textbf{Uncertainty-aware Multi-view Expert Fusion}
(\textbf{U-MoE}), a multi-view image model in which each chart view has
an expert classifier and a gate dynamically aggregates expert
predictions. Unlike ordinary MoE models that rely only on learned gate
scores, U-MoE discounts experts with high predictive entropy. This
design produces not only a final class probability, but also
interpretable quantities: expert probabilities, expert entropies, and
gate weights. These quantities allow us to ask whether the model relies
more on candlestick, moving-average, RSI, or MACD views for a given
prediction, and whether the decision is supported by confident or
uncertain experts.

Our experiments are deliberately structured around two label
definitions. The \textbf{window-trend} label indicates whether the
closing price at the end of a 30-day input window is higher than at the
start of the same window. This label is visually encoded in the chart
and measures trend recognition. The \textbf{next-day} label indicates
whether the closing price on the next trading day is higher than the
current closing price. This label measures strict future direction
prediction. Under the same chronological split, these two tasks behave
very differently. Window-trend classification reaches high accuracy,
whereas next-day prediction remains near random across CNN,
concat-fusion, and U-MoE models.

The main contributions of this paper are:

\begin{enumerate}
\def\labelenumi{\arabic{enumi}.}
\tightlist
\item
  We formulate a pure-image multi-view candlestick trend classification
  setting in which candle, moving-average, RSI, and MACD charts are
  treated as view-specific experts.
\item
  We propose U-MoE, an uncertainty-aware mixture-of-experts fusion
  method that combines expert predictions using entropy-regularized view
  weights.
\item
  We show that U-MoE provides explicit view-level interpretability
  through expert weights and predictive entropies, while retaining
  strong trend-recognition performance.
\item
  We empirically distinguish visible window-trend recognition from
  strict next-day prediction and show that U-MoE can diagnose prediction
  failure through high expert uncertainty.
\item
  We complement quantitative results with retrieval, shuffled-label
  sanity checks, frozen-probe diagnostics, and Grad-CAM visualizations.
\end{enumerate}

\subsection{2. Related Work}\label{related-work}

\subsubsection{2.1 Candlestick Image
Classification}\label{candlestick-image-classification}

Candlestick charts are widely used in technical analysis because they
summarize price action through visual primitives such as candle bodies,
shadows, gaps, and trend slopes. Several recent studies transform
financial time-series windows into chart images and apply convolutional
neural networks, object recognition methods, or image classification
architectures to identify trend or pattern classes. These methods
demonstrate that visual encodings can be useful for learning chart
structures. However, reported high accuracies often depend strongly on
label definition, window construction, and split protocol. If the target
label is derived from the same window displayed in the input image, the
task is trend recognition rather than future prediction.

\subsubsection{2.2 Multi-view Financial
Representations}\label{multi-view-financial-representations}

Financial market states can be represented from multiple perspectives:
raw price charts, moving-average overlays, momentum oscillators,
volatility indicators, and volume patterns. Multi-view learning provides
a principled way to combine such complementary observations. In this
paper, we restrict all information to rendered images, avoiding
numerical tabular features or text modalities. This allows us to study
whether a pure image model can learn and explain trend-related visual
structures from multiple technical chart views.

\subsubsection{2.3 Mixture of Experts}\label{mixture-of-experts}

Mixture-of-experts models decompose prediction into multiple specialized
experts and a gating network that assigns input-dependent expert
weights. MoE architectures are useful when different subsets of
features, tasks, or data regimes require different decision rules. In
financial chart analysis, different technical indicators may be
informative under different market conditions. This motivates a
view-level MoE design in which candlestick, MA, RSI, and MACD images are
separate experts. Compared with feature concatenation, MoE provides a
natural interpretation: the gate weight indicates how much the model
relies on each expert.

\subsubsection{2.4 Uncertainty and
Interpretability}\label{uncertainty-and-interpretability}

Uncertainty estimation is important in financial applications because a
model's abstention or lack of confidence can be as informative as its
predicted class. Predictive entropy is a simple and widely used measure
of classification uncertainty. For a two-class prediction, entropy
approaches its maximum when the model assigns nearly equal probability
to both classes. We use expert entropy to regularize view fusion and to
diagnose whether individual chart views provide reliable evidence. For
visual explanation, we also use Grad-CAM to inspect which image regions
contribute to predictions.

\subsection{3. Problem Definition}\label{problem-definition}

\subsubsection{3.1 Data Setting}\label{data-setting}

We use daily OHLC data from ten U.S. stocks: AAPL, TSLA, MSFT, AMZN,
NVDA, META, GOOG, JPM, AMD, and BAC. The historical period spans from 20
February 2020 to 18 December 2023. Each sample is constructed from a
30-trading-day window. The dataset contains 5,283 samples split
chronologically into 3,667 training samples, 818 validation samples, and
798 test samples.

The chronological split is essential. Random splits can place highly
overlapping windows from nearby dates into both training and test sets,
producing overly optimistic results. Throughout this paper, all reported
test results use a future chronological holdout.

\subsubsection{3.2 Label Definitions}\label{label-definitions}

We consider two labels.

The \textbf{window-trend} label is

{[} y\_i\^{}\{win\} = \mathbb{1}(C\_t \textgreater{} C\_\{t-29\}), {]}

where (C\_t) is the closing price at the end of the input window and
(C\_\{t-29\}) is the closing price at the start of the same 30-day
window. This label measures whether the input chart itself displays an
upward or downward trend.

The \textbf{next-day} label is

{[} y\_i\^{}\{next\} = \mathbb{1}(C\_\{t+1\} \textgreater{} C\_t). {]}

This label is outside the input window and corresponds to strict
next-trading-day direction prediction.

We treat these as separate tasks. Window-trend classification evaluates
visual trend recognition. Next-day classification evaluates whether the
image contains predictive information for the immediate future.

\subsubsection{3.3 Multi-view Image
Inputs}\label{multi-view-image-inputs}

For each window, we render four RGB image views:

\begin{itemize}
\tightlist
\item
  (x\^{}\{candle\}): raw candlestick chart.
\item
  (x\^{}\{MA\}): candlestick chart with moving-average overlays.
\item
  (x\^{}\{RSI\}): relative strength index chart.
\item
  (x\^{}\{MACD\}): MACD histogram and signal-line chart.
\end{itemize}

No numerical OHLC values, volume features, text, news, or tabular
technical indicators are provided to the model. All technical
information is visible only through images.

\subsection{4. Method}\label{method}

\subsubsection{4.1 Shared Visual Encoder}\label{shared-visual-encoder}

Let (V=\{1,\ldots,M\}) denote the set of views, with (M=4). For sample
(i), view (v) is an image (x\_i\^{}v). A shared ResNet18 encoder
(f\_\theta) maps each view to a feature vector:

{[} h\_i\^{}v = f\_\theta(x\_i\^{}v) \in \mathbb{R}\^{}d. {]}

The encoder is initialized from ImageNet weights and can optionally be
initialized from a multi-view consistency pretraining checkpoint. The
same encoder is shared across views to encourage a common visual
representation space.

\subsubsection{4.2 View-specific Experts}\label{view-specific-experts}

Each view has an expert classifier (E\_v). The expert produces logits
and probabilities:

{[} z\_i\^{}v = W\_v h\_i\^{}v + b\_v, {]}

{[} p\_i\^{}v = \mathrm{softmax}(z\_i\^{}v). {]}

The expert prediction (p\_i\^{}v) can be interpreted as the class belief
of a specific technical chart view.

\subsubsection{4.3 Expert Uncertainty}\label{expert-uncertainty}

For each expert, we compute predictive entropy:

{[} H\_i\^{}v = - \sum\_\{c=1\}\^{}\{K\}
p\_i\textsuperscript{v(c)\log p\_i}v(c), {]}

where (K=2) for binary trend classification. Entropy is high when an
expert assigns similar probabilities to upward and downward classes. In
binary classification, the maximum entropy is (\log 2 \approx 0.693).

\subsubsection{4.4 Uncertainty-aware
Gating}\label{uncertainty-aware-gating}

The gate produces an input-dependent score for each expert. In the
local-gate variant, the score is computed from each view feature:

{[} a\_i\^{}v = g\_\psi(h\_i\^{}v). {]}

In the uncertainty-aware version, the score is penalized by expert
entropy:

{[} s\_i\^{}v = a\_i\^{}v - \lambda H\_i\^{}v, {]}

where (\lambda) controls the strength of uncertainty discounting. Gate
weights are obtained by a softmax over views:

{[} \alpha\_i\^{}v = \frac{\exp(s_i^v)}{\sum_{u=1}^{M}\exp(s_i^u)}. {]}

The final U-MoE prediction is a weighted average of expert
probabilities:

{[} p\_i\^{}\{MoE\} = \sum\_\{v=1\}\^{}\{M\} \alpha\_i\^{}v p\_i\^{}v.
{]}

This formulation makes the prediction decomposable. The final
probability is determined by expert beliefs (p\_i\^{}v), expert
uncertainty (H\_i\^{}v), and gate weights (\alpha\_i\^{}v).

\subsubsection{4.5 Residual Concat Branch}\label{residual-concat-branch}

We also evaluate a residual version that combines U-MoE with a standard
concat-fusion classifier:

{[} h\_i\^{}\{cat\} = {[}h\_i\^{}1; h\_i\^{}2; \ldots; h\_i\^{}M{]}, {]}

{[} p\_i\^{}\{cat\} = \mathrm{softmax}(q\_\omega(h\_i\^{}\{cat\})). {]}

The residual prediction is

{[} p\_i = (1-\rho)p\_i\^{}\{MoE\} + \rho p\_i\^{}\{cat\}. {]}

This variant tests whether a black-box high-capacity concat branch can
improve performance while preserving expert weights for interpretation.
Empirically, however, high residual weights can cause the gate to
collapse to a single dominant view, so we treat this variant as an
ablation rather than the main method.

\subsubsection{4.6 Training Objective}\label{training-objective}

The main classification loss is the negative log-likelihood of the fused
prediction:

{[} \mathcal{L}\_\{cls\} = - \log p\_i(y\_i). {]}

We also include an auxiliary expert loss:

{[} \mathcal{L}\emph{\{aux\} = \frac{1}{M}\sum}\{v=1\}\^{}\{M\}
CE(z\_i\^{}v, y\_i), {]}

which encourages each expert to remain individually predictive. A
consistency term aligns expert distributions with the fused prediction:

{[} \mathcal{L}\emph{\{cons\} = \frac{1}{M}\sum}\{v=1\}\^{}\{M\} KL(p\_i
\parallel p\_i\^{}v). {]}

The total loss is

{[} \mathcal{L} = \mathcal{L}\emph{\{cls\} + \beta\mathcal{L}}\{aux\} +
\eta\mathcal{L}\emph{\{cons\} + \gamma\mathcal{L}}\{div\}. {]}

(\mathcal{L}\_\{div\}) is an optional feature-diversity regularizer. We
use early stopping on validation loss.

\subsection{5. Experimental Setup}\label{experimental-setup}

\subsubsection{5.1 Baselines}\label{baselines}

We compare U-MoE with:

\begin{itemize}
\tightlist
\item
  \textbf{Lightweight CNN:} a reproduction baseline for single
  candlestick-image classification.
\item
  \textbf{Mean fusion:} average of four view embeddings followed by a
  classifier.
\item
  \textbf{Concat fusion:} concatenation of four view embeddings followed
  by an MLP classifier.
\item
  \textbf{Regularized concat:} concat fusion with stronger dropout, data
  augmentation, label smoothing, and weight decay.
\item
  \textbf{Frozen encoder probe:} encoder frozen, classifier trained on
  downstream labels.
\item
  \textbf{Shuffled-label sanity check:} training labels randomly
  permuted while validation and test labels remain true.
\end{itemize}

Concat fusion is the strongest accuracy baseline but is not our proposed
contribution. It is included to test whether interpretability comes at a
performance cost.

\subsubsection{5.2 Metrics}\label{metrics}

We report accuracy, weighted F1, confusion matrices, and class-level
reports. For representation quality, we report retrieval Hit@1 using
cosine similarity between learned features. For interpretability, we
analyze expert gate weights, expert entropies, and Grad-CAM heatmaps.

\subsection{6. Results}\label{results}

\subsubsection{6.1 Window-trend Recognition vs Next-day
Prediction}\label{window-trend-recognition-vs-next-day-prediction}

Table 1 shows the effect of label definition.

\textbf{Table 1. Reproduction under different target definitions.}

\begin{longtable}[]{@{}
  >{\raggedright\arraybackslash}p{(\linewidth - 8\tabcolsep) * \real{0.1765}}
  >{\raggedright\arraybackslash}p{(\linewidth - 8\tabcolsep) * \real{0.1765}}
  >{\raggedright\arraybackslash}p{(\linewidth - 8\tabcolsep) * \real{0.1765}}
  >{\raggedleft\arraybackslash}p{(\linewidth - 8\tabcolsep) * \real{0.2353}}
  >{\raggedleft\arraybackslash}p{(\linewidth - 8\tabcolsep) * \real{0.2353}}@{}}
\toprule\noalign{}
\begin{minipage}[b]{\linewidth}\raggedright
Task
\end{minipage} & \begin{minipage}[b]{\linewidth}\raggedright
Label
\end{minipage} & \begin{minipage}[b]{\linewidth}\raggedright
Model
\end{minipage} & \begin{minipage}[b]{\linewidth}\raggedleft
Test Accuracy
\end{minipage} & \begin{minipage}[b]{\linewidth}\raggedleft
Weighted F1
\end{minipage} \\
\midrule\noalign{}
\endhead
\bottomrule\noalign{}
\endlastfoot
Visible trend recognition & window-trend & Lightweight CNN & 0.9837 &
approximately 0.9837 \\
Strict future prediction & next-day & Lightweight CNN & 0.4962 &
approximately 0.4757 \\
\end{longtable}

The same image-based setting yields very different results depending on
the target. The high window-trend score indicates that the model can
recognize the trend already displayed in the image. The next-day score
is near random, suggesting that static chart images alone do not provide
reliable immediate future-direction signals in this setting.

\subsubsection{6.2 Multi-view Fusion
Baselines}\label{multi-view-fusion-baselines}

Table 2 compares multi-view fusion strategies on the window-trend task.

\textbf{Table 2. Multi-view window-trend classification.}

\begin{longtable}[]{@{}lrrr@{}}
\toprule\noalign{}
Model & Test Accuracy & Weighted F1 & Retrieval Hit@1 \\
\midrule\noalign{}
\endhead
\bottomrule\noalign{}
\endlastfoot
Mean fusion & 0.9561 & 0.9562 & 0.9574 \\
Concat fusion & 0.9812 & 0.9813 & 0.9674 \\
Regularized concat & 0.9712 & 0.9712 & 0.9599 \\
Frozen encoder probe & 0.9398 & 0.9400 & approximately 0.9273 \\
\end{longtable}

Concat fusion achieves the best raw accuracy, confirming that
view-specific information is complementary. However, concat fusion does
not explicitly reveal which view contributes to a decision. U-MoE
addresses this interpretability gap.

\subsubsection{6.3 U-MoE on Window-trend
Classification}\label{u-moe-on-window-trend-classification}

Table 3 presents U-MoE results on the window-trend task.

\textbf{Table 3. U-MoE window-trend results.}

\begin{longtable}[]{@{}
  >{\raggedright\arraybackslash}p{(\linewidth - 8\tabcolsep) * \real{0.1667}}
  >{\raggedleft\arraybackslash}p{(\linewidth - 8\tabcolsep) * \real{0.2222}}
  >{\raggedleft\arraybackslash}p{(\linewidth - 8\tabcolsep) * \real{0.2222}}
  >{\raggedleft\arraybackslash}p{(\linewidth - 8\tabcolsep) * \real{0.2222}}
  >{\raggedright\arraybackslash}p{(\linewidth - 8\tabcolsep) * \real{0.1667}}@{}}
\toprule\noalign{}
\begin{minipage}[b]{\linewidth}\raggedright
Model
\end{minipage} & \begin{minipage}[b]{\linewidth}\raggedleft
Test Accuracy
\end{minipage} & \begin{minipage}[b]{\linewidth}\raggedleft
Weighted F1
\end{minipage} & \begin{minipage}[b]{\linewidth}\raggedleft
Hit@1
\end{minipage} & \begin{minipage}[b]{\linewidth}\raggedright
Mean Gate Weights
\end{minipage} \\
\midrule\noalign{}
\endhead
\bottomrule\noalign{}
\endlastfoot
Vanilla local MoE & 0.9674 & 0.9675 & 0.9712 & candle 0.302, MA 0.294,
RSI 0.248, MACD 0.156 \\
Uncertainty local MoE & 0.9674 & 0.9675 & 0.9712 & candle 0.291, MA
0.289, RSI 0.261, MACD 0.159 \\
Global U-MoE regularized & 0.9574 & 0.9573 & 0.9674 & candle 0.167, MA
0.386, RSI 0.035, MACD 0.412 \\
Residual MLP U-MoE, (\rho=0.5) & 0.9649 & 0.9648 & 0.9662 & candle
0.002, MA 0.635, RSI 0.352, MACD 0.011 \\
\end{longtable}

The local uncertainty-aware U-MoE achieves 96.74\% accuracy with
balanced expert weights. Although it does not exceed concat fusion in
accuracy, it provides a direct decomposition of model behavior. The
model assigns substantial weight to candlestick, MA, and RSI views,
while MACD receives lower average weight. This suggests that, for
visible trend recognition, price structure and trend/momentum views are
more consistently informative than MACD in the current setup.

The global and residual variants show that more flexible gates can
become less stable. Some variants collapse to a dominant view, such as
RSI or MA. This gate collapse is an important diagnostic: maximizing
accuracy alone may reduce interpretability and view diversity.

\subsubsection{6.4 U-MoE on Next-day
Prediction}\label{u-moe-on-next-day-prediction}

Table 4 evaluates whether U-MoE improves strict next-day prediction.

\textbf{Table 4. Next-day prediction results.}

\begin{longtable}[]{@{}
  >{\raggedright\arraybackslash}p{(\linewidth - 8\tabcolsep) * \real{0.1667}}
  >{\raggedright\arraybackslash}p{(\linewidth - 8\tabcolsep) * \real{0.1667}}
  >{\raggedleft\arraybackslash}p{(\linewidth - 8\tabcolsep) * \real{0.2222}}
  >{\raggedleft\arraybackslash}p{(\linewidth - 8\tabcolsep) * \real{0.2222}}
  >{\raggedleft\arraybackslash}p{(\linewidth - 8\tabcolsep) * \real{0.2222}}@{}}
\toprule\noalign{}
\begin{minipage}[b]{\linewidth}\raggedright
Model
\end{minipage} & \begin{minipage}[b]{\linewidth}\raggedright
Label
\end{minipage} & \begin{minipage}[b]{\linewidth}\raggedleft
Test Accuracy
\end{minipage} & \begin{minipage}[b]{\linewidth}\raggedleft
Weighted F1
\end{minipage} & \begin{minipage}[b]{\linewidth}\raggedleft
Retrieval Hit@1
\end{minipage} \\
\midrule\noalign{}
\endhead
\bottomrule\noalign{}
\endlastfoot
Lightweight CNN & next-day & 0.4962 & approximately 0.4757 & not
reported \\
Multi-view pretraining + classifier & future-k=1 & 0.4850 & 0.4300 &
0.4975 \\
U-MoE local uncertainty & future-k=1 & 0.4724 & 0.4732 & 0.4987 \\
Residual MLP U-MoE & future-k=1 & 0.4712 & 0.4305 & 0.4987 \\
\end{longtable}

U-MoE does not improve next-day prediction. This negative result is
important. A more expressive interpretable fusion model does not
overcome the lack of stable predictive signal in the image-only input.

The uncertainty estimates further support this interpretation. For local
U-MoE, mean expert entropies are:

\begin{longtable}[]{@{}lr@{}}
\toprule\noalign{}
Expert & Mean Entropy \\
\midrule\noalign{}
\endhead
\bottomrule\noalign{}
\endlastfoot
candle & 0.6789 \\
MA & 0.6797 \\
RSI & 0.6734 \\
MACD & 0.6825 \\
\end{longtable}

Since the maximum binary entropy is approximately 0.693, all experts are
highly uncertain. Thus, U-MoE explains the failure mode: the model is
not merely choosing the wrong view; rather, all views provide weak and
uncertain evidence for next-day direction.

\subsubsection{6.5 Shuffled-label Sanity
Check}\label{shuffled-label-sanity-check}

Table 5 reports the shuffled-label diagnostic.

\textbf{Table 5. Shuffled-label sanity check.}

\begin{longtable}[]{@{}
  >{\raggedright\arraybackslash}p{(\linewidth - 6\tabcolsep) * \real{0.2143}}
  >{\raggedleft\arraybackslash}p{(\linewidth - 6\tabcolsep) * \real{0.2857}}
  >{\raggedleft\arraybackslash}p{(\linewidth - 6\tabcolsep) * \real{0.2857}}
  >{\raggedright\arraybackslash}p{(\linewidth - 6\tabcolsep) * \real{0.2143}}@{}}
\toprule\noalign{}
\begin{minipage}[b]{\linewidth}\raggedright
Setting
\end{minipage} & \begin{minipage}[b]{\linewidth}\raggedleft
Test Accuracy
\end{minipage} & \begin{minipage}[b]{\linewidth}\raggedleft
Weighted F1
\end{minipage} & \begin{minipage}[b]{\linewidth}\raggedright
Confusion Matrix
\end{minipage} \\
\midrule\noalign{}
\endhead
\bottomrule\noalign{}
\endlastfoot
Concat with shuffled training labels & 0.6992 & 0.6516 & {[}{[}76,
212{]}, {[}28, 482{]}{]} \\
\end{longtable}

The performance drops substantially from 98.12\% to 69.92\%. The result
is still above 50\% because the test set is class-imbalanced and
predictions collapse toward the majority class. This diagnostic suggests
that the high window-trend result is not solely caused by file paths,
image layout, or trivial split artifacts, but it also warns that
accuracy alone is insufficient; class-level metrics and confusion
matrices are necessary.

\subsection{7. Interpretability
Analysis}\label{interpretability-analysis}

\subsubsection{7.1 Expert Weights as View-level
Explanations}\label{expert-weights-as-view-level-explanations}

U-MoE provides a view-level explanation for every prediction. For a
given sample, the final probability can be decomposed into:

{[} p(y\textbar x) = \sum\_v \alpha\_v p\_v(y\textbar x\^{}v). {]}

Here, (\alpha\_v) answers ``which chart view did the model rely on?''
and (p\_v) answers ``what did that view predict?'' This is more
transparent than concat fusion, where view information is mixed inside a
single MLP.

On window-trend classification, local U-MoE produces balanced average
weights:

{[} \alpha\emph{\{candle\}=0.291,\quad \alpha}\{MA\}=0.289,\quad
\alpha\emph{\{RSI\}=0.261,\quad \alpha}\{MACD\}=0.159. {]}

This pattern suggests that the model uses several complementary views
rather than relying only on one indicator.

\subsubsection{7.2 Entropy as Reliability
Explanation}\label{entropy-as-reliability-explanation}

Expert entropy provides a reliability signal. If an expert produces high
entropy, its prediction is uncertain and its gate score is discounted.
This is especially useful in next-day prediction. Although accuracy is
poor, the model's uncertainty is meaningful: all experts approach
maximum entropy, indicating that the model does not find reliable
view-specific evidence.

This gives U-MoE a distinctive interpretability role. It explains both
confident success on visible trend recognition and uncertainty-driven
failure on future prediction.

\subsubsection{7.3 Grad-CAM Analysis}\label{grad-cam-analysis}

We use Grad-CAM to visualize discriminative regions for each image view.
In successful window-trend cases, useful explanations should focus on
candle bodies, trend slopes, moving-average structure, RSI turning
regions, or MACD transitions rather than blank areas or image borders.

Grad-CAM is not treated as standalone proof of financial reasoning.
Instead, it complements expert weights and entropy:

\begin{itemize}
\tightlist
\item
  Gate weights identify which view is important.
\item
  Entropy indicates whether the view is confident.
\item
  Grad-CAM shows where the view-specific evidence appears in the image.
\end{itemize}

Together, these explanations form a multi-level interpretation: view
selection, view reliability, and spatial evidence.

\subsubsection{7.4 Retrieval as Representation
Evidence}\label{retrieval-as-representation-evidence}

Retrieval evaluates whether learned features organize visually similar
trend structures. On window-trend classification, U-MoE and concat
models obtain high Hit@1, around 0.96-0.97. On next-day prediction,
Hit@1 is near 0.50. This contrast supports our central claim: image
representations capture visible chart similarity, but similar chart
shapes do not reliably imply the same next-day return direction.

\subsection{8. Discussion}\label{discussion}

\subsubsection{8.1 Accuracy vs
Interpretability}\label{accuracy-vs-interpretability}

The strongest raw classifier is concat fusion, not U-MoE. This is
expected: feature concatenation preserves all view embeddings and gives
the classifier maximum flexibility. However, concat fusion behaves as a
black-box feature-level fusion method. U-MoE trades a small amount of
window-trend accuracy for explicit view-level interpretability.

Therefore, the proposed method should not be positioned as a universal
accuracy improvement over concat fusion. Its contribution is different:
it exposes which technical image experts support the decision and
whether those experts are confident.

\subsubsection{8.2 Interpreting Failure Is Also
Valuable}\label{interpreting-failure-is-also-valuable}

In financial machine learning, knowing when a model lacks reliable
evidence is critical. The next-day experiments show that U-MoE does not
improve predictive accuracy, but it reveals high expert uncertainty.
This is a useful result for responsible financial AI. An interpretable
model should not only explain confident decisions; it should also
indicate when prediction is unreliable.

\subsubsection{8.3 Implications for Candlestick Image
Research}\label{implications-for-candlestick-image-research}

Our results suggest that future candlestick-image research should
clearly separate:

\begin{enumerate}
\def\labelenumi{\arabic{enumi}.}
\tightlist
\item
  Recognizing visible trend states inside a chart.
\item
  Predicting out-of-window future returns.
\item
  Explaining which chart views drive a prediction.
\end{enumerate}

Without this separation, high classification accuracy can be
misinterpreted as market predictability. U-MoE helps make this
separation explicit by providing expert-level and uncertainty-level
diagnostics.

\subsection{9. Limitations}\label{limitations}

This study has several limitations. First, the dataset covers ten U.S.
stocks and one historical period. Broader evaluation across markets,
sectors, and regimes is needed. Second, we focus on next-day prediction
as the strict future task. Longer horizons may align better with
technical patterns and should be explored. Third, technical indicator
rendering choices may affect model behavior. Fourth, entropy is a simple
uncertainty measure; future work could investigate calibration, Bayesian
approximations, conformal prediction, or abstention mechanisms. Finally,
Grad-CAM provides qualitative evidence but does not establish causal
market reasoning.

\subsection{10. Conclusion}\label{conclusion}

This paper proposes U-MoE, an uncertainty-aware multi-view expert fusion
framework for interpretable candlestick image trend classification. By
treating candlestick, MA, RSI, and MACD charts as view-specific experts,
U-MoE produces not only a final prediction but also expert
probabilities, expert entropies, and gate weights. These quantities make
the prediction process more transparent than standard concat fusion.

Experiments show that candlestick images can support highly accurate
visible trend recognition, but strict next-day prediction remains near
random. U-MoE achieves strong window-trend performance and provides
interpretable expert weights. More importantly, when applied to next-day
prediction, U-MoE reveals uniformly high expert uncertainty, explaining
why the model cannot make reliable image-only future predictions.

The main message is therefore cautious and interpretable: candlestick
image models should not be judged solely by accuracy. A useful financial
image model should also explain which visual views support its decision
and when the evidence is too uncertain to trust. U-MoE offers a step
toward such interpretable and uncertainty-aware financial chart
analysis.

\subsection{References}\label{references}

References should be filled according to the target venue style.
Recommended citation groups include:

\begin{enumerate}
\def\labelenumi{\arabic{enumi}.}
\tightlist
\item
  Candlestick chart image classification and stock chart pattern
  recognition.
\item
  Image-based time-series trend classification.
\item
  Mixture-of-experts models and sparse/dynamic expert routing.
\item
  Multi-view representation learning.
\item
  Predictive entropy, uncertainty estimation, and model calibration.
\item
  Grad-CAM and visual explanation methods.
\item
  Financial machine learning evaluation, chronological splits, and
  leakage-aware backtesting.
\end{enumerate}

\end{document}
